\chapter{Introduzione}
\label{ch:introduzione}

La gestione della catena del freddo rappresenta una delle sfide più critiche nel settore della logistica. Il trasporto di merci deperibili richiede il mantenimento rigoroso di condizioni ambientali controllate lungo l'intera filiera distributiva, dalla produzione alla consegna finale. Anche una breve deviazione dalla temperatura richiesta può compromettere irrimediabilmente la qualità e la sicurezza dei prodotti trasportati, con conseguenze economiche rilevanti e potenziali rischi per la salute pubblica. 

A rendere il problema ancora più complesso contribuisce la frammentazione del mercato del trasporto refrigerato: da un lato, le aziende di trasporto gestiscono flotte di veicoli che spesso viaggiano con capacità parzialmente inutilizzata; dall'altro, i clienti che necessitano di spedire piccoli quantitativi si trovano costretti o a sostenere costi elevati per soluzioni FTL (\textit{Full Truckload} ovvero il noleggio di un intero veicolo refrigerato) o a ricorrere a soluzioni LTL (\textit{Less Than Truckload}, ovvero spedizioni che condividono il veicolo di trasporto con altri carichi e che vengono consolidate in centri di smistamento intermedi), che comportano molteplici passaggi in magazzino, tempi di consegna più lunghi e un rischio maggiore di rottura della catena del freddo. 

Le soluzioni attualmente disponibili sul mercato affrontano questi problemi in modo parziale e disgiunto. Piattaforme come Flock Freight \cite{flockfreight2024} hanno introdotto il concetto di \textit{Shared Truckload} (STL), sfruttando algoritmi di intelligenza artificiale per consolidare spedizioni di clienti diversi su uno stesso veicolo e ottimizzare le rotte in maniera automatica. A differenza delle spedizioni LTL tradizionali, il carico viaggia sullo stesso veicolo lungo l’intero percorso, evitando passaggi intermedi nei centri di smistamento. 

Sul versante del monitoraggio, le soluzioni IoT per la \textit{cold chain} come Lynx Fleet \cite{carrierlynx2024} operano tipicamente come sistemi B2B, scollegati dalla fase commerciale di prenotazione e pricing, e si limitano per lo più a un monitoraggio basato su soglie statiche che, pur essendo semplice da implementare, non è in grado di rilevare pattern anomali complessi dell'impianto di refrigerazione.

Il presente lavoro propone \textbf{ChillChain}, una piattaforma IoT che integra in un'unica soluzione il matching tra domanda e offerta di spedizioni refrigerate e il monitoraggio intelligente della catena del freddo con rilevazione automatica delle anomalie. L'idea alla base del progetto è consentire ai clienti di visualizzare i viaggi già programmati compatibili con le proprie esigenze di spedizione e scegliere a quale viaggio aggregare la propria spedizione. La rotta del veicolo viene ricalcolata dinamicamente per includere le nuove fermate, mentre un sistema di pricing basato sull'occupazione del mezzo applica un divisore esponenziale che rende il costo per singolo cliente progressivamente più conveniente all'aumentare delle spedizioni condivise, incentivando così un utilizzo più efficiente della capacità di trasporto.

Accanto alla componente di marketplace, la piattaforma integra un sistema di monitoraggio in tempo reale dei parametri operativi dei veicoli refrigerati. I dati telemetrici vengono trasmessi via MQTT e analizzati da un servizio dedicato di rilevazione anomalie. Quest'ultimo impiega un LSTM Autoencoder, addestrato sul comportamento nominale dell'impianto, per calcolare l'errore di ricostruzione di ciascuna finestra temporale di osservazione: quando tale errore supera una soglia appresa in fase di training, il sistema genera una notifica che viene inoltrata al tecnico responsabile attraverso la dashboard della piattaforma.

Dal punto di vista architetturale, ChillChain adotta un'architettura a microservizi composta da un backend Spring Boot che gestisce la logica applicativa e le API REST, un frontend React per l'interfaccia utente, un broker MQTT Mosquitto per la comunicazione in tempo reale, due servizi Python dedicati rispettivamente allo streaming dei dati telemetrici e alla rilevazione delle anomalie, e un database MongoDB per la persistenza. L'intera infrastruttura è orchestrata tramite Docker Compose, garantendo riproducibilità e semplicità di deployment. L'obiettivo del progetto è dimostrare la fattibilità di un approccio integrato che unisca la dimensione commerciale della logistica collaborativa con quella tecnica del monitoraggio.